\chapter{Maintaining and extending the kit}\label{ch:maintenance}

\section{Updating the kit in your governance repository}
The kit is copied, not vendored, so updates are deliberate:
\begin{enumerate}
\item Read \texttt{CHANGELOG.md} of the new version.
\item Copy the new \texttt{terraform/*.tf}, \texttt{terraform/tests/}, \texttt{scripts/} and CI
  template over yours; do \emph{not} overwrite \texttt{terraform/model/}.
\item Run \texttt{scripts/watchover test}, then open a merge request. The plan of a kit update
  should be empty; anything else needs an explanation (usually a new optional attribute).
\end{enumerate}

\section{Upgrading the Terraform provider}
The HCL relies on documented provider facts (Chapter~\ref{doc-terraform}, ``Provider facts''),
checked against v2.0.5. Before raising the version constraint:
re-check each fact in the new provider's schema and changelog, run the tests against the new
version, and review a plan against a non-production account if you have one. Watch especially for
newly ForceNew attributes and for changed defaults.

\section{Versioning and releases}
The kit and the skill share one version (\texttt{VERSION} in \texttt{scripts/watchover.py}),
recorded in \texttt{CHANGELOG.md}. CI (GitHub Actions, \texttt{.github/workflows/test.yml}) runs
the Terraform suites on both HashiCorp Terraform and OpenTofu, then the Python tests and the
simulator, then packages the skill as an artefact.

\section{Adding a kit command}
\begin{enumerate}
\item Write \texttt{scripts/<name>.py} with \texttt{build\_parser()} and \texttt{main(argv)};
  every option needs a \texttt{help} text (a test enforces it).
\item Register it in \texttt{TOOLS} in \texttt{scripts/watchover.py} with a summary and usage patterns.
\item Add tests to \texttt{scripts/tests/test\_scripts.py}.
\item Regenerate the usage guide (\texttt{scripts/watchover help --markdown > docs/usage-guide.md}).
\item Add an example execution to \texttt{examples()} in \texttt{docs/manual/build.py} (the build
  fails when a command has none) and rebuild the manual.
\end{enumerate}
Adding a skill command: add it to \texttt{SKILL\_COMMANDS} and to the command table in
\texttt{skill/akeyless-watchover/SKILL.md} (a test checks they match), and write its example session
in \texttt{docs/manual/sessions/<name>.tex}. MCP clients see the change automatically: tools and
prompts are generated from the same registry.

\section{Keeping the skill agent-neutral}
\begin{itemize}
\item Use only Agent Skills fields in the frontmatter; \texttt{validate-skill} (and CI) enforce it.
  Keep \texttt{metadata.version} equal to \texttt{VERSION}.
\item Do not name a product or a vendor tool in \texttt{SKILL.md}; a test rejects product names in its body.
  Describe capabilities instead (``if you can run a shell'').
\item Product-specific settings go into optional sidecars that other agents ignore
  (\texttt{agents/openai.yaml}).
\item When an agent changes where it looks for skills, update \texttt{AGENTS} in
  \texttt{scripts/akl\_skill.py} and the table in \texttt{docs/agents.md}.
\end{itemize}

\section{Adding a guard}
Add a precondition to \texttt{terraform\_data.model\_guards} in \texttt{terraform/main.tf}, a
negative fixture under \texttt{terraform/tests/fixtures/<name>/}, a \texttt{run} with
\texttt{expect\_failures} in \texttt{tests/model.tftest.hcl}, and the expected message to
\texttt{tests/guard\_messages.sh} so the test proves the \emph{intended} guard fired. Add the
message to the troubleshooting table (Chapter~\ref{ch:troubleshooting}).

\section{Maintaining this manual}\label{sec:manual-maintenance}
\begin{center}\small
\begin{tabularx}{\linewidth}{@{}>{\raggedright\arraybackslash}p{0.38\linewidth}>{\raggedright\arraybackslash}X@{}}
\toprule
\textbf{Content} & \textbf{Source (edit here)} \\ \midrule
Introduction, installation, quick start, workflows, security, troubleshooting, maintenance, appendices & \texttt{docs/manual/chapters/*.tex} \\
Skill command example sessions & \texttt{docs/manual/sessions/*.tex} \\
Command reference (all options, usage patterns) & \texttt{scripts/watchover.py} and each script's \texttt{build\_parser()} \\
Example executions & \texttt{examples()} in \texttt{docs/manual/build.py} (executed at build time) \\
Terraform variables and outputs & \texttt{variable}/\texttt{output} blocks in \texttt{terraform/} and \texttt{bootstrap/} \\
Concept, reference and operations chapters & \texttt{docs/*.md} (converted with pandoc) \\
Layout, title page, chapter order & \texttt{docs/manual/manual.tex} \\
\bottomrule
\end{tabularx}
\end{center}
\texttt{python3 docs/manual/build.py} regenerates \texttt{docs/manual/generated/}, executes the
examples and writes \texttt{docs/akeyless-watchover-manual.pdf}. A Python test runs
\texttt{build.py --check} (when pandoc is installed) so the generated reference parts cannot fall
behind the code; rebuild the PDF whenever you change commands, options, variables or documents.
